\chapter{Topic: Coupling, Cohesion and MVC}

\begin{aside}
Material from the remainder of the term, not covered in the original lecture
notes.
\end{aside}

\section{Measures of Design Quality}
How do we tell whether a modular design is good? Two classic measures.

\subsection{Coupling}
\begin{definition}
\textbf{Coupling} is how much distinct modules (classes) depend on each other.
\end{definition}
From low (good) to high (bad):
\begin{enumerate}
  \item modules communicate via function calls with basic parameters/results;
  \item modules pass arrays/structs back and forth;
  \item modules affect each other's control flow (e.g.\ passing flags that tell
        the other module what to do);
  \item modules share global data;
  \item modules have access to each other's implementation (\code{friend}s,
        public fields).
\end{enumerate}
With high coupling, changes to one module require changes to others, and modules
are hard to reuse individually.

\subsection{Cohesion}
\begin{definition}
\textbf{Cohesion} is how closely the elements of a single module are related to
each other.
\end{definition}
From low (bad) to high (good):
\begin{enumerate}
  \item arbitrary grouping of unrelated elements (e.g.\ \code{<utility>}:
        \code{swap}, \code{pair}, \code{move}, \dots);
  \item elements share a common theme but are otherwise unrelated, maybe sharing
        some base code (e.g.\ \code{<algorithm>});
  \item elements manipulate the same state over the lifetime of an object (e.g.\
        open/read/close a file);
  \item elements pass data to each other (a pipeline);
  \item elements cooperate to perform exactly one task.
\end{enumerate}
With low cohesion, modules are poorly organized and hard to understand and
maintain.

\begin{important}[Goal]
Low coupling, high cohesion.
\end{important}

\section{Decoupling the Interface}
Your primary program classes should not be printing things. Consider
\begin{lstlisting}[style=cpp]
class ChessBoard {
    ...
    void move(...) {
        ...
        std::cout << "Your move" << std::endl;
    }
};
\end{lstlisting}
This is bad design: it inhibits code reuse. What if we want a graphical display,
or to use the same game logic over the network, or to write output to a file?
Slightly better: give the class a stream to print to,
\begin{lstlisting}[style=cpp]
class ChessBoard {
    std::ostream &out;
  public:
    ChessBoard(std::ostream &out) : out{out} {}
    ...
};
\end{lstlisting}
but it still assumes a text interface.

\begin{definition}
\textbf{Single Responsibility Principle}: a class should have only one reason to
change.
\end{definition}
If two distinct parts of the specification can change independently (the rules
of chess vs.\ how the board is displayed), they belong in different classes.
\code{ChessBoard} should communicate with its user via parameters, results, and
possibly exceptions; a separate class handles I/O.

\section{Model--View--Controller (MVC)}
\begin{definition}
\textbf{MVC} separates a program into three parts:
\begin{itemize}
  \item the \textbf{model}: the main data and logic you are manipulating (e.g.\
        the game state);
  \item the \textbf{view}: how the model is displayed;
  \item the \textbf{controller}: how the model is manipulated (interprets user
        input, decides what to do).
\end{itemize}
\end{definition}

\begin{center}
\begin{tikzpicture}[cls/.style={draw, minimum width=24mm, minimum height=9mm},
  >={Stealth}]
\node[cls] (m) {Model};
\node[cls, below left=15mm and 6mm of m] (v) {View};
\node[cls, below right=15mm and 6mm of m] (c) {Controller};
\draw[->] (c) -- node[right, font=\small] {updates} (m);
\draw[->, dashed] (m) -- node[left, font=\small] {notifies} (v);
\draw[->] (v) -- node[below, font=\small] {user input} (c);
\end{tikzpicture}
\end{center}

\begin{itemize}
  \item The model may have multiple views (text and graphics at the same time).
        It doesn't need to know their details: the model is typically a
        \textbf{subject} and the views are \textbf{observers} (Observer pattern,
        Lecture 18).
  \item The controller mediates control flow between model and view. The model
        may (or may not) hold a reference to the controller.
  \item Swapping out the view or controller does not require changing the model,
        and the model can be tested without any UI.
\end{itemize}

\begin{lstlisting}[style=cpp]
class Game : public Subject {          // model
    Board board;
  public:
    bool makeMove(const Move &m);      // returns false if illegal
    const Board &getState() const { return board; }
};

class TextDisplay : public Observer {  // view
    Game &game;
    std::ostream &out;
  public:
    TextDisplay(Game &g, std::ostream &out) : game{g}, out{out} {
        game.attach(this);
    }
    void notify() override { out << game.getState(); }
};

class Controller {                     // controller
    Game &game;
  public:
    explicit Controller(Game &g) : game{g} {}
    void play(std::istream &in) {
        Move m;
        while (in >> m) {
            if (game.makeMove(m)) game.notifyObservers();
            else std::cerr << "Illegal move" << std::endl;
        }
    }
};
\end{lstlisting}

\section{Other Design Principles}
\begin{itemize}
  \item \textbf{Open/Closed Principle}: classes should be open for extension
        but closed for modification. Changes in a program should happen by
        writing \emph{new} code (new subclasses, new decorators, new visitors),
        not by editing working code. Strategies, Decorator, Template Method and
        Visitor all support this.
  \item \textbf{Liskov Substitution Principle}: if B is a subtype of A, a B must
        be usable anywhere an A is expected, without the client noticing.
        Subclasses must not strengthen preconditions, weaken postconditions or
        break invariants of the superclass. (Classic violation: \code{Square} as
        a subclass of a mutable \code{Rectangle} --- setting the width of a
        square also changes its height, which no Rectangle client expects.)
  \item \textbf{Program to an interface} (Lecture 17), and prefer composition
        over inheritance when the relationship is not really is-a.
\end{itemize}

\section*{Summary}
\begin{itemize}
  \item Aim for low coupling (modules depend on each other as little as
        possible) and high cohesion (each module does one coherent thing).
  \item Keep I/O out of core logic (Single Responsibility).
  \item MVC: model (data/logic, a Subject), views (Observers that display it),
        controller (turns input into model operations).
  \item Open/Closed and Liskov Substitution guide good use of inheritance.
\end{itemize}
